% GENERATED FILE. Edit canonical lessons, then run npm run generate:books.
% canonical-source-hash: fnv1a64:a3b25162a83c9ab4
% canonical-lessons: ES-C687-refugio, ES-C687-caminata, ES-C687-senda, ES-C687-trepar, ES-C687-explorar, ES-R687-repaso-680-687

\chapter{Caminos de montaña}
\label{ch:es-caminos-de-montana}
\hlchaptermodality{687}

\begin{chapteropening}
\textbf{By the end of this chapter:} \emph{I can plan a day out on foot: a long walk along a path, a mountain hut, climbing and exploring.}

Complete the last lesson of chapter 687: I can plan a day out on foot: a long walk along a path, a mountain hut, climbing and exploring.
\end{chapteropening}

\section[el refugio]{el refugio — the mountain hut; the shelter}

\label{lesson:ES-C687-refugio}



\begin{quote}
Before the new one: say the Spanish for the rock, then the Spanish for to climb.
\end{quote}

\subsection*{You'll want to know: el refugio}
\textbf{el refugio} — \textquotedblleft{}the mountain hut; the shelter\textquotedblright{}.

\begin{quote}
Dormimos en un refugio de montaña.
\end{quote}

\begin{quote}
\emph{We slept in a mountain hut.}
\end{quote}

A \emph{refugio} is a simple hut for walkers.

\begin{cousinweb}
\emph{refugio} is a learned borrowing from Latin \emph{refugium} \textquotedblleft{}a taking refuge\textquotedblright{} (Lewis \& Short), from \emph{fugere} \textquotedblleft{}to flee\textquotedblright{}.
\end{cousinweb}

\begin{grammarlens}[title={What you've built}]
One more word for telling, explaining and describing at B1.
\end{grammarlens}

\subsection*{Your turn}
\begin{itemize}
\raggedright
  \item \emph{Say it:} \emph{el refugio}
  \item \emph{Say it:} \emph{el refugio}, once more
  \item \emph{Say it:} a sentence of your own with \emph{el refugio}
  \item \emph{Recall:} say the Spanish for the rock, then the Spanish for to climb, then say \emph{el refugio} again
\end{itemize}

\begin{tcolorbox}[breakable,colback=teal!4,colframe=teal!35!black,title={Before you move on}]
What does el refugio mean? (\textquotedblleft{}The mountain hut; the shelter\textquotedblright{}.) Say it once more.
\end{tcolorbox}
\section[la caminata]{la caminata — the long walk, the hike}

\label{lesson:ES-C687-caminata}



\begin{quote}
Before the new one: say the Spanish for to climb, then the Spanish for the mountain hut.
\end{quote}

\subsection*{You'll want to know: la caminata}
\textbf{la caminata} — \textquotedblleft{}the long walk, the hike\textquotedblright{}.

\begin{quote}
Fue una caminata de cinco horas.
\end{quote}

\begin{quote}
\emph{It was a five-hour hike.}
\end{quote}

\emph{Hacer una caminata} — to go for a long walk.

\begin{cousinweb}
\emph{caminata} is built on \emph{caminar} \textquotedblleft{}to walk\textquotedblright{}, from \emph{camino} \textquotedblleft{}road\textquotedblright{}, which came into Late Latin from Celtic.
\end{cousinweb}

\begin{grammarlens}[title={What you've built}]
One more word for telling, explaining and describing at B1.
\end{grammarlens}

\subsection*{Your turn}
\begin{itemize}
\raggedright
  \item \emph{Say it:} \emph{la caminata}
  \item \emph{Say it:} \emph{la caminata}, once more
  \item \emph{Say it:} a sentence of your own with \emph{la caminata}
  \item \emph{Recall:} say the Spanish for to climb, then the Spanish for the mountain hut, then say \emph{la caminata} again
\end{itemize}

\begin{tcolorbox}[breakable,colback=teal!4,colframe=teal!35!black,title={Before you move on}]
What does la caminata mean? (\textquotedblleft{}The long walk, the hike\textquotedblright{}.) Say it once more.
\end{tcolorbox}
\section[la senda]{la senda — the path, the trail}

\label{lesson:ES-C687-senda}



\begin{quote}
Before the new one: say the Spanish for the mountain hut, then the Spanish for the long walk.
\end{quote}

\subsection*{You'll want to know: la senda}
\textbf{la senda} — \textquotedblleft{}the path, the trail\textquotedblright{}.

\begin{quote}
La senda sube por el bosque.
\end{quote}

\begin{quote}
\emph{The path goes up through the woods.}
\end{quote}

Close to \emph{el sendero}.

\begin{cousinweb}
\emph{senda} is inherited from Latin \emph{sēmita} \textquotedblleft{}a narrow way, a footpath\textquotedblright{} (Lewis \& Short; EtymDB).
\end{cousinweb}

\begin{grammarlens}[title={What you've built}]
One more word for telling, explaining and describing at B1.
\end{grammarlens}

\subsection*{Your turn}
\begin{itemize}
\raggedright
  \item \emph{Say it:} \emph{la senda}
  \item \emph{Say it:} \emph{la senda}, once more
  \item \emph{Say it:} a sentence of your own with \emph{la senda}
  \item \emph{Recall:} say the Spanish for the mountain hut, then the Spanish for the long walk, then say \emph{la senda} again
\end{itemize}

\begin{tcolorbox}[breakable,colback=teal!4,colframe=teal!35!black,title={Before you move on}]
What does la senda mean? (\textquotedblleft{}The path, the trail\textquotedblright{}.) Say it once more.
\end{tcolorbox}
\section[trepar]{trepar — to clamber, to climb (using your hands)}

\label{lesson:ES-C687-trepar}



\begin{quote}
Before the new one: say the Spanish for the long walk, then the Spanish for the path.
\end{quote}

\subsection*{You'll want to know: trepar}
\textbf{trepar} — \textquotedblleft{}to clamber, to climb (using your hands)\textquotedblright{}.

\begin{quote}
Los niños trepan a los árboles.
\end{quote}

\begin{quote}
\emph{The children climb trees.}
\end{quote}

\emph{Escalar} is the sport; \emph{trepar} is any scrambling climb.

\begin{cousinweb}
The origin of \emph{trepar} is not settled; it may imitate the sound of feet stamping.
\end{cousinweb}

\begin{grammarlens}[title={What you've built}]
One more word for telling, explaining and describing at B1.
\end{grammarlens}

\subsection*{Your turn}
\begin{itemize}
\raggedright
  \item \emph{Say it:} \emph{trepar}
  \item \emph{Say it:} \emph{trepar}, once more
  \item \emph{Say it:} a sentence of your own with \emph{trepar}
  \item \emph{Recall:} say the Spanish for the long walk, then the Spanish for the path, then say \emph{trepar} again
\end{itemize}

\begin{tcolorbox}[breakable,colback=teal!4,colframe=teal!35!black,title={Before you move on}]
What does trepar mean? (\textquotedblleft{}To clamber, to climb (using your hands)\textquotedblright{}.) Say it once more.
\end{tcolorbox}
\section[explorar]{explorar — to explore}

\label{lesson:ES-C687-explorar}



\begin{quote}
Before the new one: say the Spanish for the path, then the Spanish for to clamber.
\end{quote}

\subsection*{You'll want to know: explorar}
\textbf{explorar} — \textquotedblleft{}to explore\textquotedblright{}.

\begin{quote}
Queremos explorar la isla.
\end{quote}

\begin{quote}
\emph{We want to explore the island.}
\end{quote}

\emph{El explorador} is the explorer.

\begin{cousinweb}
\emph{explorar} is a learned borrowing from Latin \emph{explōrāre} \textquotedblleft{}to search out, to investigate\textquotedblright{} (EtymDB).
\end{cousinweb}

\begin{grammarlens}[title={What you've built}]
One more word for telling, explaining and describing at B1.
\end{grammarlens}

\subsection*{Your turn}
\begin{itemize}
\raggedright
  \item \emph{Say it:} \emph{explorar}
  \item \emph{Say it:} \emph{explorar}, once more
  \item \emph{Say it:} a sentence of your own with \emph{explorar}
  \item \emph{Recall:} say the Spanish for the path, then the Spanish for to clamber, then say \emph{explorar} again
\end{itemize}

\begin{tcolorbox}[breakable,colback=teal!4,colframe=teal!35!black,title={Before you move on}]
What does explorar mean? (\textquotedblleft{}To explore\textquotedblright{}.) Say it once more.
\end{tcolorbox}
\section[(dialogue)]{Repaso de los capítulos 680–687}

\label{lesson:ES-R687-repaso-680-687}



\begin{quote}
Say the Spanish for the cake shop and the stationer’s.
\end{quote}

\subsection*{Your turn}
Say these aloud:

\begin{itemize}
\raggedright
  \item the cake shop, the stationer’s, the hairdresser’s, the launderette, the ice-cream shop
  \item the street market, the shopping centre, the department store, holiday, to serve
  \item the street lamp, the litter bin, the large bin, the kerb, pedestrian
  \item the region, the district, the municipality, the hamlet, the mayor’s office
  \item the outsider, urban, rural, to reside, the neighbourhood
  \item the landscape, nature, the mountain range, the mountain, the desert
  \item the top, the peak, the climb, the rock, to climb
  \item the mountain hut, the long walk, the path, to clamber, to explore
\end{itemize}

\begin{tcolorbox}[breakable,colback=teal!4,colframe=teal!35!black,title={Before you move on}]
The cake shop? (\textbf{la pastelería}.) The stationer’s? (\textbf{la papelería}.) The hairdresser’s? (\textbf{la peluquería}.) The launderette? (\textbf{la lavandería}.) The ice-cream shop? (\textbf{la heladería}.) The street market? (\textbf{el mercadillo}.) The shopping centre? (\textbf{el centro comercial}.) The department store? (\textbf{los grandes almacenes}.) Holiday? (\textbf{festivo}.) To serve? (\textbf{atender}.) The street lamp? (\textbf{la farola}.) The litter bin? (\textbf{la papelera}.) The large bin? (\textbf{el contenedor}.) The kerb? (\textbf{el bordillo}.) Pedestrian? (\textbf{peatonal}.) The region? (\textbf{la región}.) The district? (\textbf{la comarca}.) The municipality? (\textbf{el municipio}.) The hamlet? (\textbf{la aldea}.) The mayor’s office? (\textbf{la alcaldía}.) The outsider? (\textbf{el forastero}.) Urban? (\textbf{urbano}.) Rural? (\textbf{rural}.) To reside? (\textbf{residir}.) The neighbourhood? (\textbf{el vecindario}.) The landscape? (\textbf{el paisaje}.) Nature? (\textbf{la naturaleza}.) The mountain range? (\textbf{la sierra}.) The mountain? (\textbf{el monte}.) The desert? (\textbf{el desierto}.) The top? (\textbf{la cima}.) The peak? (\textbf{el pico}.) The climb? (\textbf{la subida}.) The rock? (\textbf{la roca}.) To climb? (\textbf{escalar}.) The mountain hut? (\textbf{el refugio}.) The long walk? (\textbf{la caminata}.) The path? (\textbf{la senda}.) To clamber? (\textbf{trepar}.) To explore? (\textbf{explorar}.)
\end{tcolorbox}
